% GENERATED FILE. Edit canonical lessons, then run npm run generate:books.
% canonical-source-hash: fnv1a64:96d3eff57e5ad32b
% canonical-lessons: MR-C141-kinara, MR-C141-bet, MR-C141-vij, MR-C141-vadal, MR-C141-kavla

\chapter{Shore, Island, Lightning, Storm, Crow}
\label{ch:mr-shore-island-lightning-storm-crow}
\hlchaptermodality{141}

\begin{chapteropening}
\textbf{By the end of this chapter:} \emph{I can say a shore, an island, lightning, a storm and a crow.}

Complete the last lesson of chapter 141: I can say a shore, an island, lightning, a storm and a crow.
\end{chapteropening}

\section[\texorpdfstring{\mr{किनारा} (kinārā)}{किनारा (kinārā)}]{\mr{किनारा} (kinārā) — a shore, a bank}

\label{lesson:MR-C141-kinara}



\begin{quote}
Before the new one: say the Marathi for a rock, then the Marathi for a wave.
\end{quote}

\subsection*{You'll want to know: \mr{किनारा}}
\textbf{\mr{किनारा}} — \emph{kinārā} — \textquotedblleft{}a shore, a bank\textquotedblright{}.

\begin{cousinweb}
\textbf{\mr{किनारा}} (\emph{kinārā}) is a loan, usually traced to Persian \emph{kinār}, \textquotedblleft{}edge, shore\textquotedblright{}.
\end{cousinweb}

\begin{grammarlens}[title={What you've built}]
One more word for this chapter.
\end{grammarlens}

\subsection*{Your turn}
\begin{itemize}
\raggedright
  \item \emph{Say it:} \emph{kinārā}
  \item \emph{Say it:} \emph{kinārā}, once more
  \item \emph{Say it:} say \emph{kinārā}
  \item \emph{Recall:} say the Marathi for a rock, then the Marathi for a wave, then say \emph{kinārā} again
\end{itemize}

\begin{tcolorbox}[breakable,colback=teal!4,colframe=teal!35!black,title={Before you move on}]
What does \mr{किनारा} mean? (\textquotedblleft{}A shore, a bank\textquotedblright{}.) Say it once more.
\end{tcolorbox}
\section[\texorpdfstring{\mr{बेट} (beṭ)}{बेट (beṭ)}]{\mr{बेट} (beṭ) — an island}

\label{lesson:MR-C141-bet}



\begin{quote}
Before the new one: say the Marathi for a wave, then the Marathi for a shore.
\end{quote}

\subsection*{You'll want to know: \mr{बेट}}
\textbf{\mr{बेट}} — \emph{beṭ} — \textquotedblleft{}an island\textquotedblright{}.

\begin{cousinweb}
The origin of \textbf{\mr{बेट}} (\emph{beṭ}) is not settled. Neither CDIAL nor DEDR, searched by form and by meaning, has an entry for it.
\end{cousinweb}

\begin{grammarlens}[title={What you've built}]
One more word for this chapter.
\end{grammarlens}

\subsection*{Your turn}
\begin{itemize}
\raggedright
  \item \emph{Say it:} \emph{beṭ}
  \item \emph{Say it:} \emph{beṭ}, once more
  \item \emph{Say it:} say \emph{beṭ}
  \item \emph{Recall:} say the Marathi for a wave, then the Marathi for a shore, then say \emph{beṭ} again
\end{itemize}

\begin{tcolorbox}[breakable,colback=teal!4,colframe=teal!35!black,title={Before you move on}]
What does \mr{बेट} mean? (\textquotedblleft{}An island\textquotedblright{}.) Say it once more.
\end{tcolorbox}
\section[\texorpdfstring{\mr{वीज} (vīj)}{वीज (vīj)}]{\mr{वीज} (vīj) — lightning; electricity}

\label{lesson:MR-C141-vij}



\begin{quote}
Before the new one: say the Marathi for a shore, then the Marathi for an island.
\end{quote}

\subsection*{You'll want to know: \mr{वीज}}
\textbf{\mr{वीज}} — \emph{vīj} — \textquotedblleft{}lightning; electricity\textquotedblright{}.

\begin{cousinweb}
\textbf{\mr{वीज}} (\emph{vīj}) is a tadbhava, a word inherited from Sanskrit and reshaped by sound change. It comes from \emph{vidyut}, \textquotedblleft{}flashing; lightning\textquotedblright{} (CDIAL 11742).
\end{cousinweb}

\begin{grammarlens}[title={What you've built}]
One more word for this chapter.
\end{grammarlens}

\subsection*{Your turn}
\begin{itemize}
\raggedright
  \item \emph{Say it:} \emph{vīj}
  \item \emph{Say it:} \emph{vīj}, once more
  \item \emph{Say it:} say \emph{vīj}
  \item \emph{Recall:} say the Marathi for a shore, then the Marathi for an island, then say \emph{vīj} again
\end{itemize}

\begin{tcolorbox}[breakable,colback=teal!4,colframe=teal!35!black,title={Before you move on}]
What does \mr{वीज} mean? (\textquotedblleft{}Lightning; electricity\textquotedblright{}.) Say it once more.
\end{tcolorbox}
\section[\texorpdfstring{\mr{वादळ} (vādaḷ)}{वादळ (vādaḷ)}]{\mr{वादळ} (vādaḷ) — a storm}

\label{lesson:MR-C141-vadal}



\begin{quote}
Before the new one: say the Marathi for an island, then the Marathi for lightning; electricity.
\end{quote}

\subsection*{You'll want to know: \mr{वादळ}}
\textbf{\mr{वादळ}} — \emph{vādaḷ} — \textquotedblleft{}a storm\textquotedblright{}.

\begin{cousinweb}
\textbf{\mr{वादळ}} (\emph{vādaḷ}) is a tadbhava, a word inherited from Sanskrit and reshaped by sound change. It comes from \emph{vārdala}, \textquotedblleft{}rainy day, bad weather\textquotedblright{} (CDIAL 11567). CDIAL glosses the Marathi word \textquotedblleft{}cloud, bad weather\textquotedblright{}.
\end{cousinweb}

\begin{grammarlens}[title={What you've built}]
One more word for this chapter.
\end{grammarlens}

\subsection*{Your turn}
\begin{itemize}
\raggedright
  \item \emph{Say it:} \emph{vādaḷ}
  \item \emph{Say it:} \emph{vādaḷ}, once more
  \item \emph{Say it:} say \emph{vādaḷ}
  \item \emph{Recall:} say the Marathi for an island, then the Marathi for lightning; electricity, then say \emph{vādaḷ} again
\end{itemize}

\begin{tcolorbox}[breakable,colback=teal!4,colframe=teal!35!black,title={Before you move on}]
What does \mr{वादळ} mean? (\textquotedblleft{}A storm\textquotedblright{}.) Say it once more.
\end{tcolorbox}
\section[\texorpdfstring{\mr{कावळा} (kāvḷā)}{कावळा (kāvḷā)}]{\mr{कावळा} (kāvḷā) — a crow}

\label{lesson:MR-C141-kavla}



\begin{quote}
Before the new one: say the Marathi for lightning; electricity, then the Marathi for a storm.
\end{quote}

\subsection*{You'll want to know: \mr{कावळा}}
\textbf{\mr{कावळा}} — \emph{kāvḷā} — \textquotedblleft{}a crow\textquotedblright{}.

\begin{cousinweb}
\textbf{\mr{कावळा}} (\emph{kāvḷā}) is a tadbhava of Sanskrit \emph{kāka}, \textquotedblleft{}crow\textquotedblright{}, with an added \emph{-la-} (CDIAL 2993). CDIAL notes that the word imitates the bird's call.
\end{cousinweb}

\begin{grammarlens}[title={What you've built}]
One more word for this chapter.
\end{grammarlens}

\subsection*{Your turn}
\begin{itemize}
\raggedright
  \item \emph{Say it:} \emph{kāvḷā}
  \item \emph{Say it:} \emph{kāvḷā}, once more
  \item \emph{Say it:} say \emph{kāvḷā}
  \item \emph{Recall:} say the Marathi for lightning; electricity, then the Marathi for a storm, then say \emph{kāvḷā} again
\end{itemize}

\begin{tcolorbox}[breakable,colback=teal!4,colframe=teal!35!black,title={Before you move on}]
What does \mr{कावळा} mean? (\textquotedblleft{}A crow\textquotedblright{}.) Say it once more.
\end{tcolorbox}
